\documentclass[11pt,a4paper]{article}
\usepackage[T1]{fontenc}
\usepackage[utf8]{inputenc}
\usepackage[english]{babel}
\usepackage{lmodern,amsmath,amssymb,booktabs,longtable,array}
\usepackage[margin=2.3cm]{geometry}
\usepackage{xcolor,listings,hyperref,enumitem,microtype}
\hypersetup{colorlinks=true,linkcolor=blue!50!black,urlcolor=blue!60!black}
\setlist{nosep}
\lstset{basicstyle=\ttfamily\small,frame=single,breaklines=true,
columns=fullflexible,keepspaces=true,showstringspaces=false}
\newcommand{\code}[1]{\texttt{#1}}
\title{CMPUF Project Technical Documentation\\
\large Stochastic Model, Test Data, Solution Methods, and Results}
\author{CMPUF Project}
\date{\today}
\begin{document}
\maketitle
\tableofcontents
\newpage

\section{Purpose and Software Scope}

The \code{cmpuf} program solves a capacitated facility location problem in
which capacities are subject to random disruptions. Facility locations are
first-stage decisions shared by every scenario. Once a capacity scenario is
known, customer flows are reoptimized. The implementation is therefore a
discrete two-stage stochastic program with continuous recourse.

The maintained executable combines:
\begin{itemize}
\item \code{cpp/cmpuf.cpp}: input, scenarios, indicators, and main program;
\item \code{cpp/modele.cpp}: CPLEX model, greedy heuristic, and local search;
\item \code{cpp/pmedian.cpp}: p-median preprocessing;
\item \code{cpp/transport.cpp}: independent scenario models;
\item \code{cpp/tri.cpp}: distance-based neighborhood sorting.
\end{itemize}

\section{Notation and Input Parameters}

Let $N=\{0,\ldots,n-1\}$ be both the customer set and the candidate-location
set, $I=\{0,\ldots,m-1\}$ the facility set,
$R=\{0,\ldots,ncap-1\}$ the capacity-state set, and
$\Omega=\{0,\ldots,K-1\}$ the scenario set, where $K=ncap^m$.

\begin{longtable}{>{$}l<{$}p{11cm}}
\toprule
\text{Symbol}&\text{Meaning}\\ \midrule \endhead
w_j&Demand of customer $j$.\\
H=\sum_jw_j&Total demand.\\
d_{lj}&Distance or cost from location $l$ to customer $j$, stored as
\code{profit[l][j]}.\\
D_{\max}=\max_{l,j}d_{lj}&Largest distance in the instance.\\
C_i^m&Demand assigned to facility $i$ by p-median preprocessing.\\
\gamma&Nominal-capacity multiplier.\\
S_i=\gamma C_i^m&Nominal capacity of facility $i$.\\
a_r&Probability weight read from \code{walpha}.\\
b_r&Capacity multiplier read from \code{wbeta}.\\
\alpha&Overall disruption parameter.\\
C_{ki}&Capacity of facility $i$ in scenario $k$.\\
P_k&Probability of scenario $k$.\\ \bottomrule
\end{longtable}

The member \code{beta} is read, stored, and printed, but it does not occur in
an active capacity formula. Effective multipliers are read directly from
\code{wbeta}.

\section{Capacity and Scenario Generation}

P-median capacities are scaled as
\[
S_i=\gamma C_i^m.
\]
The disrupted and nominal states are
\begin{align}
\operatorname{cap}_{ir}&=b_rS_i &&r=0,\ldots,ncap-2,\\
\operatorname{cap}_{i,ncap-1}&=S_i.
\end{align}
The supplied parameter file therefore produces $0$, $0.33S_i$, $0.66S_i$,
and $S_i$.

Marginal probabilities are
\begin{align}
p_r&=\alpha a_r &&r=0,\ldots,ncap-2,\\
p_{ncap-1}&=1-\alpha.
\end{align}
The $a_r$ values must sum to one. The supplied weights satisfy
$0.10+0.40+0.50=1$.

A scenario is $T_k=(T_{k0},\ldots,T_{k,m-1})\in R^m$. The recursive
\code{Init\_T} function enumerates the full Cartesian product and computes
\[
C_{ki}=\operatorname{cap}_{i,T_{ki}},
\qquad
P_k=\prod_{i\in I}p_{T_{ki}}.
\]
This produces $K=ncap^m$ scenarios and approximately $n^2K$ flow variables.
Scenario enumeration is the main scalability limitation.

\section{Implemented Two-Stage Stochastic Model}

\subsection{Variables}

First-stage variables are
\[
x_{il}=\begin{cases}
1&\text{if facility $i$ is located at node $l$,}\\
0&\text{otherwise.}
\end{cases}
\]
They are created in $[0,1]$ and converted to integer variables through
\code{IloConversion}. Recourse variables
$z_{ljk}\geq0$ represent the amount shipped from location $l$ to customer
$j$ in scenario $k$.

\subsection{Expected-Utility Objective}

The implemented objective is
\begin{equation}
\max\quad
\sum_{k\in\Omega}P_k\sum_{l\in N}\sum_{j\in N}
(D_{\max}-d_{lj})z_{ljk}.
\label{eq:objective}
\end{equation}
When total delivered quantity is fixed, this is equivalent to expected
transportation-cost minimization. The transformation gives nonnegative
coefficients.

\paragraph{Implementation caveat.}
Demand is imposed as a lower bound rather than an equality. With excess
capacity and positive $D_{\max}-d_{lj}$, the solver may ship more than demand.
In that case, \eqref{eq:objective} is not strictly equivalent to minimizing
cost for exactly satisfied demand.

\subsection{Constraints}

Each facility is located exactly once:
\begin{equation}
\sum_{l\in N}x_{il}=1,\qquad i\in I.
\label{eq:location}
\end{equation}
In the active path, these constraints are added progressively by the greedy
heuristic and remain in the model used for the exact solve.

For every location and scenario:
\begin{equation}
\sum_{j\in N}z_{ljk}\leq\sum_{i\in I}C_{ki}x_{il},
\qquad l\in N,\ k\in\Omega.
\end{equation}
For every customer and scenario:
\begin{equation}
\sum_{l\in N}z_{ljk}\geq w_j,
\qquad j\in N,\ k\in\Omega.
\end{equation}

The optional constraint $\sum_i x_{il}\leq1$ is implemented by
\code{Init\_contraintes\_x\_x}, but its call is commented out. Facility
co-location is therefore allowed.

The $x$ variables are here-and-now decisions, while $z$ variables are
scenario-dependent recourse. Benders annotations place $x$ in the master and
partition $z$ by scenario. The exact run uses CPLEX's automatic Benders
strategy ($-1$).

\section{P-Median Preprocessing}

The preprocessing model uses $X_l\in\{0,\ldots,m\}$, the number of medians at
node $l$, and continuous flows $z^P_{lj}$. It imposes
\begin{align}
\sum_lX_l&=m,\\
z^P_{lj}&\leq w_jX_l,\\
\sum_lz^P_{lj}&=w_j.
\end{align}
It maximizes $\sum_{l,j}(D_{\max}-d_{lj})z^P_{lj}$, which is equivalent to
distance minimization because demand is satisfied at equality. Since $X_l$
may exceed one, p-median co-location is allowed.

The solution is expanded into $m$ individual facilities. For each one,
$C_i^m$ is the flow assigned to its median and becomes its reference
capacity in the stochastic model.

\section{Solution Algorithms}

\subsection{P-Median Evaluation}

The p-median location vector is fixed in the stochastic model and all
scenario flows are optimized. This produces \code{VMEDIAN1}. A local search
from the same vector produces \code{VMEDIAN2}.

\subsection{Greedy Construction}

\code{CMpuf::MyGreedyHeuristic} locates facilities sequentially:
\begin{enumerate}
\item facilities $j>i$ are temporarily fixed to zero;
\item constraint \eqref{eq:location} is added for facility $i$;
\item CPLEX solves the partial model;
\item the best location of $i$ is recorded and fixed;
\item the temporary constraint is removed.
\end{enumerate}
The final value is \code{VGREEDY}.

\subsection{Local Search}

\code{Localsearch} performs best-improvement single-facility relocation. All
locations are fixed, each facility is released in turn, and CPLEX finds its
best new location. The best strictly improving move is accepted until no
improvement remains.

If the radius is smaller than $n$, each facility may move only among the
\code{rayon} closest nodes to its initial position. \code{ROPTri} sorts these
nodes. Local search is run from the greedy solution
(\code{VGREEDYLS}) and from the p-median solution (\code{VMEDIAN2}).

\subsection{Exact MIP and Scenario Models}

If argument 2 is \code{e}, CPLEX solves the full MIP with the requested time
limit and relative gap. Internal heuristic frequency and automatic Benders
strategy are set to $-1$. The code accepts \emph{Optimal},
\emph{OptimalTol}, \emph{SolLim}, and \emph{AbortTimeLim}; consequently,
\code{VOPT} may be an incumbent at the time limit rather than a proven
optimum. A best heuristic vector is computed, but its \code{MIPStart} call is
commented out.

After an exact run, \code{ResolutionScenarios} solves a separate
transportation/location model for every known-capacity scenario and fills
\code{SolScenarios} and \code{SolSupply}. Detailed scenario display is
implemented but disabled in the main program.

\section{Location Indicators}

For $L=(L_0,\ldots,L_{m-1})$, centralization is the mean pairwise facility
distance:
\[
\operatorname{cent}(L)=
\frac{\sum_{0\leq i<j<m}d_{L_iL_j}}{m(m-1)/2}.
\]
If $q(L)$ is the number of distinct occupied locations, co-location is
\[
\operatorname{coloc}(L)=100\frac{m-q(L)}m.
\]

\section{Test File Formats}

\subsection{Instance-List File}

Argument 1 is not an instance itself. It is a text file containing one
instance path per line:
\begin{lstlisting}
/absolute/path/to/instance1
/absolute/path/to/instance2
\end{lstlisting}
For example, \code{../BENCHMARK/calgaryhospital15} contains the path of the
actual instance. Paths must not contain spaces because they are read with
\code{fscanf(file, "\%s", filename)}.

\subsection{Standard or Hospital Instance}

When \code{INPUT\_TYPE} is not \code{geometric}, the file contains:
\begin{lstlisting}
n
w[0] ... w[n-1]
d[0][0] ... d[0][n-1]
...
d[n-1][0] ... d[n-1][n-1]
\end{lstlisting}

\subsection{Geometric Instance}

When \code{INPUT\_TYPE} is exactly \code{geometric}, the reader expects:
\begin{lstlisting}
n
mtmp
ptmp[0] ... ptmp[mtmp-1]
w[0] ... w[n-1]
d[0][0] ... d[n-1][n-1]
\end{lstlisting}
The temporary \code{ptmp} array is read but not subsequently used. Demand and
matrix values are extracted through Concert stream operators.

\subsection{Stochastic Parameter File}

\code{parameters/param} follows:
\begin{lstlisting}
ncap
a[0] ... a[ncap-2]
b[0] ... b[ncap-2]
\end{lstlisting}
The supplied contents are:
\begin{lstlisting}
4
0.10 0.40 0.50
0.00 0.33 0.66
\end{lstlisting}

\section{Build and Test Procedure}

From the project root, compile with:
\begin{lstlisting}[language=bash]
make clean
make -j2 cmpuf
\end{lstlisting}
or:
\begin{lstlisting}[language=bash]
cmake --preset release
cmake --build --preset cmpuf --parallel
\end{lstlisting}
The executable is \code{bin/cmpuf}.

The program expects nine user arguments:
\begin{lstlisting}[language=bash]
./bin/cmpuf INSTANCE_LIST MODE M GAP LOCAL_RADIUS \
  BB_RADIUS TIME_LIMIT PARAMETER_FILE INPUT_TYPE
\end{lstlisting}

\begin{longtable}{clp{8.8cm}}
\toprule
No.&Argument&Purpose\\ \midrule \endhead
1&\code{INSTANCE\_LIST}&File containing one instance path per line.\\
2&\code{MODE}&\code{e} enables the exact solution; other values run only
heuristic evaluations.\\
3&\code{M}&Number of facilities.\\
4&\code{GAP}&CPLEX relative gap.\\
5&\code{LOCAL\_RADIUS}&Local-search neighborhood size.\\
6&\code{BB\_RADIUS}&Intended local branch-and-bound radius; calls are
currently commented out.\\
7&\code{TIME\_LIMIT}&CPLEX limit in seconds.\\
8&\code{PARAMETER\_FILE}&File containing $ncap$, $a_r$, and $b_r$.\\
9&\code{INPUT\_TYPE}&\code{geometric}, or conventionally \code{hospital} for
the standard format.\\ \bottomrule
\end{longtable}

Example:
\begin{lstlisting}[language=bash]
./bin/cmpuf ../BENCHMARK/calgaryhospital15 e 4 0.00 \
  5 40 30000 parameters/param hospital
\end{lstlisting}

\subsection{Implemented Test Plan}

For each instance, p-median preprocessing is run once. Then $\gamma$ takes
$1$, $1.25$, $1.50$, and $1.75$, while $\alpha$ ranges from $0$ to $1$ in
steps of $0.1$. The value of $\beta$ remains zero. For every pair, the
p-median location, greedy solution, and both local searches are evaluated.
Mode \code{e} additionally runs the exact MIP. This gives theoretically
$4\times11=44$ configurations per instance, subject to floating-point loop
rounding.

Recommended checks are:
\begin{enumerate}
\item run \code{./bin/cmpuf} without arguments and verify the argument-count
error;
\item begin with a small instance and short time limit;
\item verify $\sum_{r<ncap-1}a_r=1$;
\item ensure $n^2ncap^m$ fits available memory;
\item verify that \code{results/hospital} exists and is writable;
\item compare all five objective values.
\end{enumerate}

\section{Result File Formats}

\subsection{Naming}

Four hard-coded prefixes are used:
\begin{lstlisting}
results2026_v1_
solgreedy_v1_
solopt_v1_
solcent_v1_
\end{lstlisting}
The suffix includes part of \code{argv[1]} followed by
\code{\_M\_GAP\_LOCAL\_RADIUS\_BB\_RADIUS\_TIME\_LIMIT}.

\paragraph{Current limitation.}
\code{concat} uses \code{argv[1]+16} rather than extracting the final path
component. Names depend on the instance-list path length and may be truncated,
for example \code{garyhospital15}. Output directories are hard-coded absolute
paths. Every output is opened in append mode.

\subsection{Main results2026 File}

Each numeric row contains 25 ampersand-separated fields:
\begin{longtable}{clp{8.5cm}}
\toprule
No.&Field&Meaning\\ \midrule \endhead
1&\code{n}&Number of nodes.\\
2&\code{m}&Number of facilities.\\
3--5&\code{Alpha}, \code{Beta}, \code{Gamma}&Test parameters.\\
6&\code{VOPT}&Exact value, or $-1$ outside exact mode.\\
7&\code{VGREEDY}&Greedy value.\\
8&\code{VGREEDYLS}&Local-search value from greedy.\\
9&\code{VMEDIAN1}&Stochastic value of p-median locations.\\
10&\code{VMEDIAN2}&Local-search value from p-median.\\
11--15&\code{cent\_*}&Centralization indicators.\\
16--20&\code{coloc\_*}&Co-location percentages.\\
21&\code{topt}&Exact-solution CPLEX time.\\
22&\code{tgreedy}&Greedy time.\\
23&\code{tgreedyls}&Local-search time from greedy.\\
24&\code{tpmedian1}&P-median preprocessing time.\\
25&\code{tpmedian2}&Local-search time from p-median.\\ \bottomrule
\end{longtable}

\paragraph{Current header anomaly.}
The header call has 24 \code{\%s} placeholders for 25 labels and no line
break. Thus \code{tpmedian2} is omitted and the first data row is joined to
the header. Later numeric rows still contain the documented 25 fields.

\subsection{solgreedy, solopt, and solcent}

For every $(\alpha,\gamma)$ pair, \code{solgreedy} receives:
\begin{lstlisting}
Lls1[0] Lls1[1] ... Lls1[m-1]
\end{lstlisting}
This is the local-search solution from greedy, not the raw greedy vector.

In exact mode, \code{solopt} receives:
\begin{lstlisting}
Lopt[0] Lopt[1] ... Lopt[m-1]
\end{lstlisting}

Also in exact mode, each tab-separated \code{solcent} row is:
\begin{lstlisting}
alpha    beta    gamma    VOPT    centralizationopt
\end{lstlisting}

Standard output prints a LaTeX-table-style row for every configuration and
can be redirected to a report file.

\section{Interpretation Caveats}

\begin{itemize}
\item The objective is expected utility, and its cost interpretation requires
fixed delivered quantities.
\item Demand uses $\geq$, so mathematical over-service is possible.
\item Facility co-location is allowed.
\item \code{beta} currently has no effect.
\item scenarios and variables grow exponentially with $m$.
\item \code{VOPT} at the time limit may not be proven optimal.
\item repeated runs append to existing output files.
\item \code{while(!feof(file1))} may attempt an additional iteration.
\end{itemize}

\section{Code Traceability}

\begin{longtable}{p{5.2cm}p{9cm}}
\toprule
Element&Main function or file\\ \midrule \endhead
Instance and parameter loading&\code{CMpuf} constructor,
\code{cpp/cmpuf.cpp}\\
Capacities&\code{CMpuf::Init\_cap}\\
Scenarios&\code{CMpuf::Init\_T}, \code{CMpuf::Init\_C}\\
Variables&\code{CMpuf::Init\_x}, \code{CMpuf::Init\_z}\\
Constraints&\code{CMpuf::Init\_contraintes},
\code{CMpuf::Init\_contraintes\_x}\\
Objective&\code{CMpuf::Init\_obj}\\
P-median&\code{cpp/pmedian.cpp}\\
Greedy heuristic&\code{CMpuf::MyGreedyHeuristic}\\
Local search&\code{CMpuf::Localsearch}\\
Exact solution&\code{CMpuf::Resolution} and the main program\\
Scenario models&\code{CMpuf::ResolutionScenarios},
\code{cpp/transport.cpp}\\
Output&\code{AffichageSolutions}, \code{AffichageSolutionsScenarios},
\code{AffichageCoutTransport}\\ \bottomrule
\end{longtable}

\end{document}
